\section*{अध्याय \ref{ch:g9:inside-the-atom} --- परमाणु के भीतर}

\begin{solution}{exo:g9:inside-the-atom:1}
\omterm{def:g9:inside-the-atom:nucleus}{नाभिक} में: \omterm{def:g9:inside-the-atom:proton}{प्रोटॉन} (धनात्मक) और \omterm{def:g9:inside-the-atom:proton}{न्यूट्रॉन} (कोई आवेश नहीं)। उसके चारों ओर:
\omterm{def:g9:inside-the-atom:nucleus}{इलेक्ट्रॉन} (ऋणात्मक)।
\end{solution}

\begin{solution}{exo:g9:inside-the-atom:2}
\ce{^{16}_{8}O}: 8 \omterm{def:g9:inside-the-atom:proton}{प्रोटॉन}, 8 \omterm{def:g9:inside-the-atom:proton}{न्यूट्रॉन}, 8 \omterm{def:g9:inside-the-atom:nucleus}{इलेक्ट्रॉन}।
\ce{^{27}_{13}Al}: 13, 14, 13। \ce{^{40}_{20}Ca}: 20, 20, 20।
\end{solution}

\begin{solution}{exo:g9:inside-the-atom:3}
उसमें \omterm{def:g9:inside-the-atom:nucleus}{इलेक्ट्रॉन} (हर एक पर आवेश $-e$) उतने ही हैं जितने \omterm{def:g9:inside-the-atom:proton}{प्रोटॉन} (हर एक पर
आवेश $+e$): आवेश एक-दूसरे को निरस्त कर देते हैं।
\end{solution}

\begin{solution}{exo:g9:inside-the-atom:4}
$A = 26 + 30 = 56$: \ce{^{56}_{26}Fe}।
\end{solution}

\begin{solution}{exo:g9:inside-the-atom:5}
हाँ: दोनों का $Z = 17$ है, इसलिए दोनों क्लोरीन हैं। उनकी \omterm{def:g9:inside-the-atom:atomic-number}{द्रव्यमान संख्याएँ}
अलग हैं (18 और 20 \omterm{def:g9:inside-the-atom:proton}{न्यूट्रॉन}): वे \omterm{def:g9:inside-the-atom:isotope}{समस्थानिक} हैं।
\end{solution}

\begin{solution}{exo:g9:inside-the-atom:6}
नहीं: उनके \omterm{def:g9:inside-the-atom:atomic-number}{परमाणु क्रमांक} अलग हैं (6 और 7), इसलिए वे अलग-अलग तत्व हैं,
कार्बन और नाइट्रोजन। \omterm{def:g9:inside-the-atom:isotope}{समस्थानिकों} का $Z$ एक ही होता है।
\end{solution}

\begin{solution}{exo:g9:inside-the-atom:7}
$56 \times 1.67 \times 10^{-27} \approx \qty{9.35e-26}{kg}$।
\end{solution}

\begin{solution}{exo:g9:inside-the-atom:8}
\omterm{def:g7:atoms-and-molecules:atom}{परमाणु} ज़्यादातर ख़ाली जगह है: \omterm{def:g9:inside-the-atom:nucleus}{नाभिक} बहुत छोटा है, इसलिए ज़्यादातर कणों के
रास्ते में कोई \omterm{def:g9:inside-the-atom:nucleus}{नाभिक} नहीं आता और वे सीधे आर-पार निकल जाते हैं।
\end{solution}

\begin{solution}{exo:g9:inside-the-atom:9}
$10^{-10} / 10^{-15} = 10^{5}$: एक लाख गुना छोटा।
\end{solution}

\begin{solution}{exo:g9:inside-the-atom:10}
रसायन \omterm{def:g9:inside-the-atom:nucleus}{इलेक्ट्रॉनों} पर निर्भर करता है, और किसी तत्व के \omterm{def:g9:inside-the-atom:isotope}{समस्थानिकों} में
\omterm{def:g9:inside-the-atom:nucleus}{इलेक्ट्रॉनों} की संख्या एक ही होती है (वही $Z$)।
\end{solution}

\begin{solution}{exo:g9:inside-the-atom:11}
$26 \times 9.11 \times 10^{-31} = \qty{2.37e-29}{kg}$। भिन्न:
$2.37 \times 10^{-29} / 9.35 \times 10^{-26} \approx 2.5 \times
10^{-4}$, यानी द्रव्यमान का लगभग \qty{0.025}{\percent}।
\end{solution}

\begin{solution}{exo:g9:inside-the-atom:12}
ताँबा-63 के 6915 \omterm{def:g7:atoms-and-molecules:atom}{परमाणु} और ताँबा-65 के 3085 \omterm{def:g7:atoms-and-molecules:atom}{परमाणु}। औसत:
$(6915 \times 63 + 3085 \times 65)/10\,000 = 63.6$ \omterm{def:g9:inside-the-atom:proton}{न्यूक्लिऑन}।
\end{solution}

\begin{solution}{pb:g9:inside-the-atom:1}
\textbf{1.} \ce{^{35}_{17}Cl} और \ce{^{37}_{17}Cl}।

\textbf{2.} क्लोरीन-35: 17 \omterm{def:g9:inside-the-atom:proton}{प्रोटॉन}, 18 \omterm{def:g9:inside-the-atom:proton}{न्यूट्रॉन}, 17 \omterm{def:g9:inside-the-atom:nucleus}{इलेक्ट्रॉन}।
क्लोरीन-37: 17 \omterm{def:g9:inside-the-atom:proton}{प्रोटॉन}, 20 \omterm{def:g9:inside-the-atom:proton}{न्यूट्रॉन}, 17 \omterm{def:g9:inside-the-atom:nucleus}{इलेक्ट्रॉन}।

\textbf{3.} दोनों में 17 \omterm{def:g9:inside-the-atom:proton}{प्रोटॉन} हैं, इसलिए दोनों क्लोरीन तत्व हैं; वे केवल
\omterm{def:g9:inside-the-atom:proton}{न्यूट्रॉनों} की संख्या में भिन्न हैं, इसलिए वे \omterm{def:g9:inside-the-atom:isotope}{समस्थानिक} हैं।

\textbf{4.} हाँ: रसायन \omterm{def:g9:inside-the-atom:nucleus}{इलेक्ट्रॉनों} पर निर्भर करता है, और दोनों में 17 हैं।

\textbf{5.} $35 \times 1.67 \times 10^{-27} = \qty{5.845e-26}{kg}$।

\textbf{6.} $17 \times 9.11 \times 10^{-31} = \qty{1.55e-29}{kg}$:
\omterm{def:g9:inside-the-atom:nucleus}{नाभिक} से लगभग 4000 गुना कम, नगण्य।

\textbf{7.} $37 \times 1.67 \times 10^{-27} = \qty{6.179e-26}{kg}$।

\textbf{8.} $\qty{1}{g} = \qty{e-3}{kg}$; $10^{-3} / 5.845 \times
10^{-26} \approx 1.71 \times 10^{22}$ \omterm{def:g7:atoms-and-molecules:atom}{परमाणु}।

\textbf{9.} \qty{75.8}{\percent} और \qty{24.2}{\percent}।

\textbf{10.} $758 \times 5.845 \times 10^{-26} = \qty{4.430e-23}{kg}$;
$242 \times 6.179 \times 10^{-26} = \qty{1.495e-23}{kg}$।

\textbf{11.} $4.430 \times 10^{-23} + 1.495 \times 10^{-23} =
\qty{5.925e-23}{kg}$।

\textbf{12.} $5.925 \times 10^{-23} / 1000 \approx
\qty{5.93e-26}{kg}$: क्लोरीन \omterm{def:g7:atoms-and-molecules:atom}{परमाणु} का औसत द्रव्यमान।
\end{solution}
